% \subsection{Problem Motivation}
% Modern text-to-image diffusion models have achieved remarkable success in generating high-quality, photorealistic images (\citep{dhariwal2021diffusionmodelsbeatgans}, \citep{9880056}, \cite{poole2022dreamfusiontextto3dusing2d} and \citep{saharia2022photorealistictexttoimagediffusionmodels}). However, they continue to struggle with compositional understanding (\citep{chatterjee2024gettingrightimprovingspatial}, \citep{adamkiewicz2026prettyisntusefulinvestigating}), which is the ability to correctly interpret and render spatial relationships, object interactions, and logical constraints specified in natural language prompts. While models like Stable Diffusion 3.5 (SD3.5) (\citep{esser2024scalingrectifiedflowtransformers}, \citep{rombach2022highresolutionimagesynthesislatent}) can generate high-quality images of individual objects, they systematically fail when prompts require precise spatial positioning, logical negation \citep{ban2024understandingimpactnegativeprompts}, or relational reasoning. This limitation represents a critical gap between language comprehension and visual generation that impacts the practical utility of these models for tasks requiring precise spatial control.

% The research problem stems from the training nature of standard Vision-Language Models (VLMs) \citep{bordes2024introductionvisionlanguagemodeling}. Models like CLIP \citep{radford2021learningtransferablevisualmodels} are often trained on loose image-caption pairs where spatial precision is not a priority. Prior work has primarily focused on scaling diffusion architectures and improving image quality metrics Fréchet Inception Distance \citep{heusel2018ganstrainedtimescaleupdate}, Importance Sampling \citep{grover2019biascorrectionlearnedgenerative}, CLIP scores \citep{hessel2022clipscorereferencefreeevaluationmetric}, but systematic investigation of compositional failures remains limited. Most studies evaluate object presence rather than spatial relationships, masking fundamental limitations in language-to-visual grounding. Consequently, the text embeddings fed into the diffusion transformer may lack the ``spatial coordinates'' necessary for the denoising process to arrange objects correctly.

% Understanding and addressing compositional failures in diffusion models has several important implications:



% \subsubsection{Scientific Value}
% Compositionality is a fundamental property of human language and cognition. Investigating why state-of-the-art vision language models fail at seemingly simple spatial reasoning tasks, such as positioning a bed above a book, provides insights into the limitations of current training paradigms and model architectures.

% \subsubsection{Practical Value}
% Applications requiring spatial precision (\citep{ruiz2023dreamboothfinetuningtexttoimage}, \citep{tian2023stablerepsyntheticimagestexttoimage}, \citep{hammoud2024synthclipreadyfullysynthetic} and \citep{voetman2023bigdatamythusing}), are currently limited by the inability of models to reliably interpret spatial language. Improving compositional reasoning would significantly expand the utility of text-to-image systems in professional creative workflows.


% \subsection{Architectural Insight}
% By systematically replacing specific components, specifically text encoders, and measuring their impact; this work contributes to understanding which model components are responsible for compositional understanding versus general image quality. Our research focuses on whether targeting the text encoding stage can resolve these issues without requiring full model retraining.



% \subsection{Research Aim and Overview}
% In this work, we systematically quantify failure rates across various spatial and relational concepts in SD3.5. We investigate whether replacing SD3.5's CLIP-L encoder with CLIC (Compositional Language-Image Contrastive) \citep{peleg2025advancingcompositionalawarenessclip}, which is trained on compositional image-text pairs, enhances compositional reasoning.

% We detail a technical methodology for replacing the model's internal text encoder with specialized versions: CLIC RedCaps \footnote{https://huggingface.co/nmndeep/CLIC-ViT-L-14-224-PixPr-RedCaps}, trained on 12M Reddit image-text pairs rich in detailed spatial descriptions \citep{desai2021redcapswebcuratedimagetextdata}, and CLIC CogVLM\footnote{https://huggingface.co/nmndeep/CLIC-ViT-L-14-224-CogVLM}, which is aligned with high-performance visual features \citep{wang2024cogvlmvisualexpertpretrained}. We analyze performance across 13 phrases: above, below, to the left of, to the right of, between, behind, closeby, far from, in front of, on top of, under, over and without; using automated VLM evaluation via Qwen2-VL \citep{bai2023qwenvlversatilevisionlanguagemodel} and measure image quality using Fréchet Inception Distance to ensure image fidelity is preserved. Additionally, we develop automated VLM-based evaluation infrastructure to assess these relationships at scale and document the complete pipeline for converting and integrating OpenCLIP models into production diffusion pipelines.

% The remainder of this report is organized as follows: Section \ref{sec:methods} details our data collection and architectural methodology, Section \ref{sec:results} presents quantitative and qualitative results, and Section \ref{sec:conclusion} discusses findings, limitations, and future work.

\subsection{Problem Motivation}
Modern text-to-image diffusion models have achieved remarkable success in generating high-quality, photorealistic images \citep{dhariwal2021diffusionmodelsbeatgans, 9880056, poole2022dreamfusiontextto3dusing2d, saharia2022photorealistictexttoimagediffusionmodels}. However, they continue to struggle with compositional understanding \citep{chatterjee2024gettingrightimprovingspatial, adamkiewicz2026prettyisntusefulinvestigating}, defined as the ability to correctly interpret and render spatial relationships, object interactions, and logical constraints specified in natural language prompts. While models such as Stable Diffusion 3.5 (SD3.5) \citep{esser2024scalingrectifiedflowtransformers, rombach2022highresolutionimagesynthesislatent} can generate stunningly detailed images of individual objects, they systematically fail when prompts require precise spatial positioning, logical negation \citep{ban2024understandingimpactnegativeprompts}, or complex relational reasoning. This limitation represents a critical gap between language comprehension and visual generation, ultimately restricting the practical utility of these models for tasks that demand rigorous spatial control.

This research problem fundamentally stems from the nature of the training data used for standard Vision-Language Models (VLMs) \citep{bordes2024introductionvisionlanguagemodeling}. Foundational models like CLIP \citep{radford2021learningtransferablevisualmodels} are frequently trained on vast datasets of loose, web-scraped image-caption pairs where strict spatial precision is rarely a priority. Historically, prior work has primarily focused on scaling diffusion architectures and optimizing general image quality metrics, such as Fréchet Inception Distance (FID) \citep{heusel2018ganstrainedtimescaleupdate}, Importance Sampling \citep{grover2019biascorrectionlearnedgenerative}, and standard CLIP scores \citep{hessel2022clipscorereferencefreeevaluationmetric}; while the systematic investigation of compositional failures has remained largely limited. Because most evaluation studies verify mere object presence rather than analyzing accurate spatial relationships, fundamental limitations in language-to-visual grounding are often masked. Consequently, the text embeddings fed into the diffusion transformer may inherently lack the explicit ``spatial coordinates'' necessary for the denoising process to arrange generated objects correctly.

Understanding and addressing compositional failures in diffusion models has several important implications:

\subsubsection{Scientific Value}
Compositionality is a foundational property of both human language and cognitive processing. Investigating precisely why state-of-the-art vision-language models fail at seemingly simple spatial reasoning tasks—such as accurately positioning a bed above a book—provides critical insights into the inherent limitations of current training paradigms and model architectures.

\subsubsection{Practical Value}
Downstream applications requiring rigorous spatial precision \citep{ruiz2023dreamboothfinetuningtexttoimage, tian2023stablerepsyntheticimagestexttoimage, hammoud2024synthclipreadyfullysynthetic, voetman2023bigdatamythusing} are currently severely bottlenecked by the inability of underlying models to reliably interpret spatial language. Consequently, improving a model's compositional reasoning capabilities would significantly expand the reliability and utility of text-to-image systems, particularly within professional creative workflows and controlled generation pipelines.


\subsection{Architectural Insight}
By systematically isolating and replacing specific model components (namely, the text encoders) and comprehensively measuring the resulting impact, this work contributes to a deeper understanding of which architectural elements are primarily responsible for compositional understanding as opposed to general image quality. Crucially, our research focuses on determining whether specifically targeting the text encoding stage can resolve these spatial reasoning deficits without incurring the prohibitive computational costs of full model retraining.


\subsection{Research Aim and Overview}
In this work, we systematically quantify the failure rates across a diverse set of spatial and relational concepts within SD3.5. We specifically investigate whether replacing SD3.5's default CLIP-L encoder with CLIC (Compositional Language-Image Contrastive) \citep{peleg2025advancingcompositionalawarenessclip}, an encoder explicitly trained on compositionally dense image-text pairs, measurably enhances the model's overall compositional reasoning capabilities.

To this end, we detail a comprehensive technical methodology for replacing the model's internal text encoder with two specialized variants: CLIC RedCaps\footnote{\url{https://huggingface.co/nmndeep/CLIC-ViT-L-14-224-PixPr-RedCaps}}, trained on 12 million Reddit image-text pairs rich in detailed spatial descriptions \citep{desai2021redcapswebcuratedimagetextdata}, and CLIC CogVLM\footnote{\url{https://huggingface.co/nmndeep/CLIC-ViT-L-14-224-CogVLM}}, which is aligned with high-performance visual features \citep{wang2024cogvlmvisualexpertpretrained}. We analyze model performance across 13 distinct relational phrases (above, below, to the left of, to the right of, between, behind, closeby, far from, in front of, on top of, under, over, and without) by utilizing automated VLM-based evaluation via Qwen3-VL \citep{bai2023qwenvlversatilevisionlanguagemodel}. Furthermore, we measure overall image quality using the Fréchet Inception Distance to verify that image fidelity is preserved despite the encoder swap. In addition to our empirical findings, we develop an automated, VLM-based evaluation infrastructure designed to assess these spatial relationships at scale, and we document the complete pipeline required for converting and integrating OpenCLIP models into production-level diffusion pipelines.

The remainder of this report is organized as follows: Section \ref{sec:methods} details our data collection processes and architectural methodology; Section \ref{sec:results} presents both the quantitative and qualitative results of our experiments; and Section \ref{sec:conclusion} discusses our core findings, acknowledges current limitations, and outlines directions for future work.